% 04_features.tex -- the rehauled paper's section, as bullets to expand into prose (2026-10-15).
% Sources and pointers behind every number: writeup/PAPER_REHAUL_OUTLINE_2026-10-15.md, section 4.
% House rules: every number in the final prose comes from a generated macro (writeup/generated/*.tex), never
% typed; park (comment out with a dated note) rather than delete.
% Macros to use here: \bkNSeries, \campPNewBase, \campPNewAll, \campKeptBase, \campKeptAll, \campNSessionEdge
% 2026-10-15 (notation pass): symbols and wording per writeup/NOTATION_GUIDE_2026-10-15.md (bar t -> interval
% (t,i); ladder \bar u_{t,k} -> HAR-type average \bar u^{(k)}_{t,i}; activity flag alpha -> chi (alpha is the
% elastic-net mixing weight); r_{t,k} -> IQR^{*(k)}_t (r is the return); \mathcal K_t -> \mathcal J_t (\mathcal K
% is the horizon set); asinh scale s_t -> kappa_{t,i} (s^2 is the periodicity factor); T_W -> W; q_{p,t} -> Q_p).

% 2026-10-15 (consolidation): window row s -> t' (s is the periodicity factor); calendar dummies \bm d_{t,i} ->
% \bm D_{t,i} (d_t is the DM loss differential); typed counts -> macros where one exists (47 -> \rhfiNExog, 628 ->
% \cmPNewAll, 22 -> \cmPNewBase, 369 -> \cmKeptAll); "session-open/close dummies" -> opening/closing-interval dummies.
% 2026-10-15 (finish, user): "intraday periodicity adjustment" -> "diurnal adjustment" (term rename only); Table 3
% (tab:rh_dm_design) parked, so its two \ref's are dropped (old lines parked in place).

% 2026-10-15 (notation pass): parked -- section title; "features" only for ML methods, "predictors" in prose
% \section{Features}
\section{Predictors}
\label{sec:rh_predictors}
% 2026-10-17 (user): prose rewritten in the style of Zhang, Zhang, Cucuringu, and Qian (2024).
\begin{itemize}
% 2026-10-18 (user, second-list N7): parked -- series count from a non-FI macro (\bkNSeries, make_buckets_tex.py, the sum of the \bkSize* group sizes; same 47), so an FI rerun cannot change it. Was: "\item \textbf{Eight groups of exogenous predictors (\rhfiNExog{} series in total).} As listed in"
\item \textbf{Eight groups of exogenous predictors (\bkNSeries{} series in total).} As listed in
  % 2026-10-17 (user): parked -- undefined ref (the buckets table is parked). Was: "As summarized in Table~\ref{tab:rh_buckets},"
  Appendix Table~\ref{tab:rh_dm_dictionary}, the exogenous series span eight economic categories: realized moments (\bkSizeMom),
  liquidity (\bkSizeLiq), equal-weighted market moments (\bkSizeMew), value-weighted market moments (\bkSizeMvw),
  sentiment (\bkSizeSent),
  % 2026-10-15 (user): parked -- "park all tables that I haven't mentioned in this chat". Was: implied volatility (\bkSizeIv), volatility demand (\bkSizeVd), FOMC (\bkSizeFomc); listed in Appendix
  implied volatility (\bkSizeIv), volatility demand (\bkSizeVd), and FOMC announcements (\bkSizeFomc).
  Section~\ref{sec:rh_data} describes the construction of each category, and Appendix
  Table~\ref{tab:rh_dm_dictionary} provides the definition of each individual series:
  % 2026-10-15 (user): parked -- "park all tables that I haven't mentioned in this chat"
  % Table~\ref{tab:rh_dm_dictionary}.
  % 2026-10-17 (user, E5): the dictionary is live in the appendix. The pointer above replaces this comment.
  \[
    % 2026-10-18 (user, second-list N7): parked -- \rhfiNExog -> \bkNSeries (non-FI). Was: "\rhfiNExog = \bkSizeMom + \bkSizeLiq + \bkSizeMew + \bkSizeMvw + \bkSizeSent + \bkSizeIv + \bkSizeVd + \bkSizeFomc ."
    \bkNSeries = \bkSizeMom + \bkSizeLiq + \bkSizeMew + \bkSizeMvw + \bkSizeSent + \bkSizeIv + \bkSizeVd + \bkSizeFomc .
  \]
  % 2026-10-15 (finish, phase 2): typed group sizes 7, 14, 6, 6, 3, 3, 4, 4 -> \bkSize... (same values).
  \begin{itemize}
  \item Specifically, the four FOMC series are defined as
    \[
      \mathbf{1}\{\text{announcement in } (t,i)\}, \quad
      \mathbf{1}\{\text{announcement on day } t\}, \quad 1/(1 + n^{+}_{t,i}), \quad 1/(1 + n^{-}_{t,i}),
    \]
  \item where $n^{\pm}_{t,i}$ denotes the number of intervals until the next ($+$) and since the last ($-$)
    scheduled announcement.
    % loading.py:183-221: distances on the weekend-cut grid, before the dead-session drop; all four missing
    % after the last print (2023-11-01). The 2020-03-03 exclusion is a no-op (outline flag 13).
  \end{itemize}
\item \textbf{Predictors constructed from each series.}
  \begin{itemize}
  % 2026-10-17 (user, new-notes 6): parked -- unit of k stated (30-minute grid intervals, both models; har.py:48 rolls over the rows of the full bar series, lag_scope "global", segmentation.py:39-43). Was: "over horizons of $k$ intervals, lagged by one interval"
  \item Following \citet{Corsi2009}, we construct HAR-type trailing averages of each transformed series $\tilde u$
    % 2026-10-18 (user, new-notes 6): parked -- k in samples. Was: "over horizons of $k$ 30-minute grid intervals, lagged by one interval:"
    over horizons of $k$ samples (one sample is one 30-minute grid interval), lagged by one sample:
    \[
      \bar u^{(k)}_{t,i} = \frac{1}{k} \sum_{j=1}^{k} \tilde u_{t,i-j}, \qquad k \in \mathcal{K},
    \]
    % 2026-10-17 (user): parked -- closing-interval material (closing horizon set and rows i = c moved to Appendix app:rh_closing_predictors). Was:
    % where the horizon set $\mathcal{K}$ is given by
    % \[
    %   \mathcal{K} =
    %   \begin{cases}
    %     \{1, 5, 25, 125, 625, 3125\} & \text{closing-interval model},\\
    %     \{2^0, 2^1, \dots, 2^{11}\} \ \text{(base 2, \rhdmNHorizonsPooled{} horizons)} & \text{intraday model}.
    %   \end{cases}
    % \]
    % \item Here, $k$ counts retained intervals across the full intraday grid $\mathcal{G}$ (including overnight
    %   sessions), while the closing-interval specification restricts estimation to the rows $i = c$.
    % 2026-10-19 (user): parked -- base 5 across the paper. Was:
    % where the horizon set of the intraday model is
    % \[
    %   \mathcal{K} = \{2^0, 2^1, \dots, 2^{11}\} \quad \text{(base 2, \rhdmNHorizonsPooled{} horizons)}.
    % \]
    where the horizon set is
    \[
      \mathcal{K} = \{1, 5, 25, 125, 625, 3125\} \quad \text{(base 5, \rhdmNHorizonsPooled{} horizons)}.
    \]
  % 2026-10-17 (user, new-notes 6): parked -- unit of k stated (30-minute grid intervals, both models; har.py:48 rolls over the rows of the full bar series, lag_scope "global", segmentation.py:39-43). Was: "Here, $k$ counts retained intervals across the full intraday grid"
  % 2026-10-18 (user, new-notes 6): parked -- k in samples. Was: "\item Here, $k$ counts retained 30-minute intervals of the full intraday grid $\mathcal{G}$ (including overnight"
  \item Here, $k$ counts samples, the retained intervals of the full intraday grid $\mathcal{G}$ (including overnight
    % 2026-10-19 (user): parked -- base 5 across the paper; both models share the horizon set. Was: "... The horizon set of the closing-interval specification is given in Appendix~\ref{app:rh_closing_predictors}."
    sessions), in the intraday and the closing-interval model alike. The closing-interval specification uses the same horizon set
    (Appendix~\ref{app:rh_closing_predictors}).

    % har.py:18-48: rolling(k, min_periods=1).mean().shift(1). Per bar: production base-5 ladder, lag_scope
    % "global" (experiments/capture_design_close.py:116-117; segmentation.py:39-43), so k is in BARS, not
    % sessions (the previous bullet said sessions). Pooled: HAR_BASE=2 gives 1..2048 (specs/causal_tune_linear.py:
    % 198-209). sections/data_prep.tex eq:dp_har writes base 2 for the per-bar ladder: wrong, it is base 5.
    % 2026-10-19 (user): base 5 across the paper -- the pooled models are refitted at base 5 (HAR_BASE unset, the
    % production ladder); the base-2 pooled fits of record (a0, the pooled trees of 7 October) are replaced.
  \item In addition, we compute the corresponding $k$-interval trailing averages of the availability indicator,
    denoted $\bar a^{(k)}_{t,i}$.
  % 2026-10-17 (user, E6): parked -- "zero-inflated" -> "mode-inflated". Was: "... $\bar\chi^{(k)}_{t,i}$, for 7 zero-inflated series."
  % 2026-10-17 (user): parked -- closing-interval material (the count of 7 is the closing-interval panel's, note E6; moved to Appendix app:rh_closing_predictors). Was: "for 7 mode-inflated series."
  \item We also compute $k$-interval trailing averages of an activity indicator, $\bar\chi^{(k)}_{t,i}$, for
    \rhdmNActivePooled{} mode-inflated series in the intraday model.
  % 2026-10-17 (user, E6): parked -- "zero-inflated" -> "mode-inflated"; closing-panel chi-bar wording (the count of 7, \rhdmNActive, the "three counts" sentences) removed; the pile-of-zeros sentence rewritten as a mode at 30; decision line narrowed to the intraday count. Was:
  % \editc{red}{E6}{Decision: one rule and one count.\\Problem: the bullet uses 7 zero-inflated series, the transformation uses the $\rhdmNActive$ series that pass a 20\% mode rule, and the intraday column count uses $\rhdmNActivePooled$. Those are three different sets.\\
  % This bullet says the activity average $\bar\chi^{(k)}_{t,i}$ is built for 7 series, and it calls those series zero-inflated. That word means the raw series has a pile of exact zeros. The transformation further down does not look for zeros. It sets the indicator $\chi_{t,i}$ to 1 when the value differs from the series' full-sample mode, and it keeps a series only when that mode covers at least 20\% of intervals. A mode can sit at a number other than zero, so the 20\% rule and a mass at zero select different series. The number of series that pass the 20\% rule is $\rhdmNActive$. The intraday column formula then multiplies by a third count, $\rhdmNActivePooled$. The activity block of the design is the product of one of these counts with the horizons, and the three counts are not the same count.}
  \editc{red}{E6}{Decision: one rule and one count, the intraday count $\rhdmNActivePooled$.\\Problem: the bullet calls the series mode-inflated, the transformation keeps the series that pass a 20\% mode rule, and the intraday column count uses $\rhdmNActivePooled$.\\
This bullet calls these series mode-inflated. That word means the raw series has a mode at 30, and 30 is the number of observations: the count of one-minute intervals in a half-hour bar. The transformation further down does not look for zeros. It sets the indicator $\chi_{t,i}$ to 1 when the value differs from the series' full-sample mode, and it keeps a series only when that mode covers at least 20\% of intervals. A mode can sit at a number other than zero, so the 20\% rule and a mass at zero select different series. The intraday column formula multiplies by $\rhdmNActivePooled$.\\
Additional decision (2026-10-17, user): remove the trailing averages of the activity indicators, the $\bar\chi^{(k)}_{t,i}$ blocks, from the model.\\
Note (2026-10-17, user): the count of 7 comes from the closing-interval panel and belongs in the economic application, the section on the last-30-minute closing trade; it may move to the appendix, with a reference back from that section.}
  \item For notational simplicity in our tables and figures, we suppress the overbar and the interval index $i$
    % 2026-10-17 (user): parked -- closing-interval material (example horizon 25 is on the closing ladder; the intraday ladder is base 2). Was: "($i = c$ in the closing-interval specification) ... $\mathrm{VIX}^{(25)}_t$ ... over the 25"
    ($i = c$ for the scored forecast), writing $u^{(k)}_t$ for series $u$ averaged over $k$
    % 2026-10-17 (user, new-notes 6): parked -- unit of k stated (30-minute grid intervals, both models; har.py:48 rolls over the rows of the full bar series, lag_scope "global", segmentation.py:39-43). Was: "averaged over the 16 intervals ending at 15:30"
    % 2026-10-18 (user, new-notes 6): parked -- k in samples. Was: "intervals; for instance, $\mathrm{VIX}^{(16)}_t$ denotes the transformed VIX series averaged over the 16"
    % 2026-10-19 (user): parked -- base 5 across the paper (16 is not a base-5 horizon; 25 is). Was: "samples; for instance, $\mathrm{VIX}^{(16)}_t$ denotes the transformed VIX series averaged over the 16"
    % 2026-10-19 (user): series names through \ser (raw names for now). Was: "$\mathrm{VIX}^{(16)}_t$"
    samples; for instance, $\ser{vix}{\mathrm{VIX}}^{(25)}_t$ denotes the transformed VIX series averaged over the 25
    % 2026-10-18 (user, new-notes 6): parked -- k in samples. Was: "30-minute intervals ending at 15:30 on day $t$. We denote the implied volatility series by $\mathrm{VIX}$,"
    % 2026-10-19 (user): raw parquet names for now. The \else branch is the sentence of the symbols, unchanged; the
    % superscript convention does not carry over to the raw names, which end in _ewstock / _vwstock instead.
    \ifrawseries
    samples ending at 15:30 on day $t$. We refer to each exogenous series by its column name in the data
    (Table~\ref{tab:rh_dm_dictionary}): the implied volatility series are \sername{vix}, \sername{vix3m}, and
    \sername{vvix}; the suffixes \sername{_ewstock} and \sername{_vwstock} mark the equal- and value-weighted
    constituent aggregates (e.g., \sername{ofi_ewstock}, \sername{sumret2_vwstock}, and
    \sername{spread_vwstock}); and the announcement predictors are \sername{fomc_release}, \sername{fomc_day},
    \sername{fomc_until_inv}, and \sername{fomc_since_inv}, while writing
    \else
    samples ending at 15:30 on day $t$. We denote the implied volatility series by $\mathrm{VIX}$,
    $\mathrm{VIX3M}$, and $\mathrm{VVIX}$, label equal- and value-weighted constituent aggregates with the
    superscripts EW and VW (e.g., $\mathrm{OFI}^{\mathrm{EW}}$, $\RV^{\mathrm{VW}}$, and
    $\mathrm{Spread}^{\mathrm{VW}}$), and denote the announcement predictors by $\mathrm{FOMC}$, while writing
    \fi
    $\bar a^{(k)}_t(u)$ and $\bar\chi^{(k)}_t(u)$ for the trailing averages of the availability and activity
    indicators.
  \end{itemize}
\item \textbf{HAR benchmark.} The HAR benchmark predictor set consists of the HAR components of the target,
  $\bar y^{(k)}_{t,i}$—constructed via the same trailing-average operator as $\bar u^{(k)}_{t,i}$ applied to the
  % 2026-10-17 (user, new-notes 7): parked -- the intraday model has \rhdmNHorizonsPooled{} HAR columns (base-2 ladder, specs/causal_tune_linear.py:197-209). Was: "winsorized target (6 columns)"
  % 2026-10-19 (user): base 5 across the paper -- \rhdmNHorizonsPooled{} is now the base-5 count (the ladder of the closing-interval model).
  winsorized target (\rhdmNHorizonsPooled{} columns in the intraday model)—together with the calendar vector $\bm D_{t,i}$, which comprises five
  day-of-the-week dummy variables, the hour of the interval-end timestamp, and the overnight session indicator.
  % 2026-10-16 (formula pass): parked -- bar y^{(k)}_{t,i} = (1/k) sum_{j=1}^{k} tilde y_{t,i-j} repeated the
  % average already displayed for bar u^{(k)}_{t,i}.
  % 2026-10-17 (user, E9): parked -- 16 columns, the 5+4+7 split, and everything after is_overnight.
  % Was: and 16 calendar dummies D_{t,i} (day of the week, monthly option expiration, quarterly quadruple
  % witching, month end, quarter end and others).
  % \[ 16 = 5 day of week + 4 time of day + 7 expiration, month, quarter. \]
  % calendar.py:28-35 (DOW_0-4, hour, is_overnight, is_open, is_close); expiry.py:120-135 (opex, opex week,
  % quad witch, rebalance close, month end, quarter end, days_to_opex).
  \editc{orange}{E9}{Decision: the calendar stops at the overnight indicator, and that indicator stays only if the overnight intervals stay in the training sample (note E15).\\Problem: the columns after the overnight indicator are out of the specification, and the printed column counts and Table~\ref{tab:rh_oos} still include them. Whether the overnight indicator itself stays depends on the sample in note E15.\\
The sentence above keeps three calendar pieces: the five weekday dummies, the hour of the bar-end stamp, and the overnight indicator. Every column that used to sit after the overnight indicator is out of the specification. That is the opening-interval indicator, the closing-interval indicator, expiration, expiration week, quadruple witching, the quadruple-witching close, month end, quarter end, days to expiration, and the products of the HAR components with the opening and closing indicators. The overnight indicator's own place depends on which sample note E15 keeps. On all $\rhdmNIntervals$ intervals it equals 1 off the cash session and again on the 16:00 bar, so it varies across the window and a regression can estimate it. On the 13 regular-hours bars it equals 1 only on the 16:00 bar, so inside that sample it is constant on every bar but the close. The column counts in the formulas below, and the fits in Table~\ref{tab:rh_oos}, were built with the dropped columns still in the matrix.}
  % 2026-10-17 (user): parked -- the two subbullets under the HAR benchmark. The drop list stays in note E9.
  % The closing-interval screen is in Methods. Was:
  % \begin{itemize}
  % \item Dropped, everything after the overnight indicator: the opening-interval and closing-interval indicators,
  %   monthly expiration, expiration week, quadruple witching, the quadruple-witching closing auction, month end,
  %   quarter end, trading days to expiration, and the products of the target HAR components with the opening- and
  %   closing-interval indicators.
  % \item At the closing interval, hour and the overnight indicator are constant, so the column screen drops them
  %   in every closing-interval estimation.
  % \end{itemize}
  % 2026-10-15 (user): parked -- Table 3 (tab:rh_dm_design) is parked. Was:
  % drops these five at every estimation of the closing-interval model (Table~\ref{tab:rh_dm_design}).
  % 2026-10-17 (user, E9): parked -- the note that asked for the names. Was:
  % \editc{blue}{E9}{Decision: name the $\rhdmNCalendar$ dummies and write the interaction. The count splits them into 5 day-of-week, 4 time-of-day, and 7 expiration, month, and quarter. Prose names day of week, monthly expiration, quadruple witching, month end, and quarter end. This item names the five dropped at the close: hour, overnight, opening interval, closing interval, and the quadruple-witching closing-auction dummy. The other members of the 7 are still ``and others.'' The $\rhdmNEdgePooled$ columns are ``HAR components times the opening- and closing-interval dummies,'' with no product and no clock time for the opening dummy.}
  % 2026-10-15 (consolidation): added to match tab:rh_dm_design (four constant dummies) and tab:rh_fi_linear (five
  % never in a fit; addenda A.3).
\item \textbf{Sequential feature transformations.} Each exogenous series undergoes the following seven
  transformations in sequence:
  \begin{enumerate}
  \item \emph{Diurnal adjustment:} we divide each raw series $u_{t,i}$ by a time-of-day scale factor
    $b^{(u)}_{t,i}$ estimated over the same prior-day set $\mathcal{D}_{t,i}$ as $s^2_{t,i}$ and bounded
    relative to its own trailing median:
    \[
      \tilde u_{t,i} = \frac{u_{t,i}}{b^{(u)}_{t,i}}, \qquad
      b^{(u)}_{t,i} =
      \begin{cases}
        \clip\bigl(\operatorname{mean}_{t' \in \mathcal{D}_{t,i}} u_{t',i},\ 0.5\,\rho_{t,i},\ 2\rho_{t,i}\bigr),
          & u \ge 0,\\[2pt]
        \max\bigl\{\clip\bigl(\operatorname{sd}_{t' \in \mathcal{D}_{t,i}} u_{t',i},\ 0.5\,\rho_{t,i},\
          2\rho_{t,i}\bigr),\ 0.1\,\bar\sigma^{(u)}\bigr\}, & u \text{ signed},
      \end{cases}
    \]
    \begin{itemize}
    \item where $\rho_{t,i}$ denotes the median of the 250 most recent unbounded diurnal factors for interval $i$
      (requiring at least 20 observations),
    \item $\bar\sigma^{(u)}$ denotes the full-sample median of the within-interval standard deviation on the
      same intraday interval,
    \item $b^{(u)}_{t,i}$ is set to $1$ whenever the factor of a nonnegative series is zero or undefined, and
    \item no diurnal adjustment is applied to the implied volatility indices (VIX, VIX3M, and VVIX), the
      StockTwits sentiment series, or the FOMC series.
    \end{itemize}
    % target.py:186-300. The band acts on exogenous series only (target: slot_band = 0). rho_{t,i} has no extra
    % shift: each baseline in it is already lagged (target.py:201-204). Baselines are estimated on the
    % forward-filled u (ffill executor.py:540 precedes robust_transform :562).
  \item \emph{Variance-stabilizing transformation:} we apply a domain-specific mapping
    $\tilde u_{t,i} \leftarrow g(\tilde u_{t,i})$ according to the first matching rule in
    \[
      g(v) =
      \begin{cases}
        \sqrt{v} & \text{squared returns, realized variance, turnover,}\\
          & \text{bipower variation, effective spread},\\
        \sign(v)\,|v|^{1/2} & \text{autocovariance},\\
        v^{1/3} & \text{cubed returns},\\
        \sign(v)\,|v|^{1/4} & \text{fourth-power returns},\\
        v & \text{other signed series, absolute returns},\\
        \log \max(v, 10^{-12}) & \text{otherwise}.
      \end{cases}
    \]
    \begin{itemize}
    \item For signed series without a specific power-root rule whose standard-deviation divisor hits its floor on
      at least 50\% of intervals (namely, the volatility-demand series), we instead apply an inverse hyperbolic
      sine transformation directly to the raw series:
      \[
        \tilde u_{t,i} = \operatorname{asinh}(u_{t,i} / \kappa_{t,i}), \qquad
        \kappa_{t,i} = \max\bigl\{ e_{t,i-1},\ 0.01 \operatorname{med}_{j \ge 1} e_{t,i-j} \bigr\},
      \]
    % 2026-10-18 (user, second-list 18): parked -- half-life is counted in nonzero intervals (zero rows are blanked and
    % skipped, ignore_na=True), not days. Was: "where $e_{t,i}$ denotes an exponentially weighted moving average of
    % $|u|$ computed over nonzero intervals ($u \ne 0$) with a 250-day half-life."
    % 12,000 = ASINH_SCALE_HALFLIFE_DAYS (250) x PERIODS_PER_DAY (48); src/features/transforms/target.py:95,451.
    \item where $e_{t,i}$ denotes an exponentially weighted moving average of $|u|$ computed over nonzero
      intervals ($u \ne 0$) with a half-life of 12{,}000 such intervals (nominally 250 days of 48 intervals).
    \end{itemize}
    % target.py:387-442 (map), 303-345 (asinh), 644-663 (routing). asinh replaces both the diurnal division and
    % g: it acts on raw u. The routing share is full-sample. fomc_release in {0,1} maps to {log 1e-12, 0}; after
    % the 1/99 clip adj_fomc_release is constant (outline flag 13).
  % 2026-10-16 (formula pass): parked -- the clip repeated the target's winsorization (Methods), with the
  % window 240 and the 1st/99th percentiles typed a second time. Was:
  % tilde u <- clip(tilde u, Q_{0.01}(U_{t,i}), Q_{0.99}(U_{t,i})),
  % U_{t,i} = {tilde u_{t,i-j} : j = 1, ..., 240}.
  \item \emph{Winsorization:} each series is clipped at the same rolling quantile thresholds used for the target
    (see Section~\ref{sec:rh_methods_setup}), evaluated over the series' own preceding intervals.
  \item \emph{Missing-value imputation:} any missing observation is imputed using the expanding median of the
    series' previously observed transformed values, while retaining the availability indicator $a_{t,i}$ as a
    separate predictor:
    \[
      \tilde u_{t,i} \leftarrow
      \begin{cases}
        \tilde u_{t,i}, & a_{t,i} = 1,\\
        \operatorname{med}\mathcal{A}_{t,i}, & a_{t,i} = 0,\ \mathcal{A}_{t,i} \ne \emptyset,\\
        \tilde u_{(1)}, & a_{t,i} = 0,\ \mathcal{A}_{t,i} = \emptyset,
      \end{cases}
    \]
    \begin{itemize}
    \item where $\mathcal{A}_{t,i} = \{\tilde u_{t',i'} : (t',i') \prec (t,i),\ a_{t',i'} = 1\}$, with
      $(t',i') \prec (t,i)$ indicating that $(t',i')$ is a retained interval strictly prior to $(t,i)$, and
    \item $\tilde u_{(1)}$ denotes the transformed value on the series' first observed interval.
    \end{itemize}
    % 2026-10-16 (user): causal neutral fill. executor.py causal_neutral_fill, CAUSAL_NEUTRAL_FILL:
    % expanding median of observed adjusted values strictly before the row; the first observed value
    % before the series starts. Parked -- the full-sample median:
    % med{tilde u_{t',i'} : a_{t',i'} = 1}, with no restriction to earlier intervals.
    % The parked branch is adj.fillna(adj.median()); the old comment cited executor.py:602-608.
  \item \emph{Activity indicators:} for the \rhdmNActive{} series whose full-sample mode accounts for at least
    20\% of all intervals, we define a binary activity indicator
    \[
      \chi_{t,i} = \mathbf{1}\{u_{t,i} \ne \operatorname{mode}(u)\}.
    \]
    % CODE BUG (outline flag 1): u_{t,i} here is the FORWARD-FILLED raw value (executor.py:618-627 runs after the
    % ffill at :540), so chi_{t,i} = 1 on dead tails: fomc_release / fomc_day after 2023-11-01, voldemand after
    % 2023-08-31. Intended: chi_{t,i} = a_{t,i} * 1{u_{t,i} != mode}.
    % 2026-10-18 (user, second-list N4): the causal rule is the specification; this item and Table 5's chi-bar column marked rerun-needed.
    \editc{orange}{S4a}{Decision: the causal rule in src/backtest/executor.py (CAUSAL\_MODE\_INFLATION, ACTIVE\_FROM\_OBSERVED; commit b29774a, 15 October 2026) is the specification. Rerun needed under that rule.\\Problem: this item describes the old rule, a full-sample mode covering at least 20\% of intervals with the indicator read after the forward fill. Under the causal rule the mode is decided from earlier prints only and the flags come from observed prints, so the volatility-demand series are never mode-inflated and the StockTwits sentiment and message-count series are mode-inflated in some years. The printed forecasts (pooled linear models of 6 October, pooled trees of 7 October 2026) predate the rule and use the old one; so does the $\bar\chi$ column of Appendix Table~\ref{tab:rh_dm_dictionary}, whose generator is not re-run (E5). Rewrite this item and regenerate the table after the rerun.}
  \item \emph{HAR aggregation:} we form the trailing averages $\bar u^{(k)}_{t,i}$, $\bar a^{(k)}_{t,i}$,
    $\bar\chi^{(k)}_{t,i}$, and $\bar y^{(k)}_{t,i}$ across horizons $k \in \mathcal{K}$.
  \item \emph{Robust standardization:} under the pooled intraday scheme, we standardize the trailing averages of
    series values over the rolling estimation window using the median and interquartile range (with lower
    floors). Specifically, the median and interquartile range are computed exclusively over intervals with
    $a_{t',i'} = 1$ inside the rolling window $\mathcal{W}_{t,i}$ (see Appendix~\ref{sec:rh_exp_schemes}),
    re-evaluated every 480 intervals and held constant until the next grid point. Let $M^{(k)}_{t,i}$ denote the
    maximum such interquartile range observed across all grid points up to $(t,i)$. We apply a $0.01$ relative
    % 2026-10-18 (user, new-notes 8): parked -- forward reference to the IQR* definition added. Was: "floor in the denominator (note that $\mathrm{IQR}^{*(k)}$ does not enter this pooled denominator):"
    floor in the denominator (note that $\mathrm{IQR}^{*(k)}$, defined in Appendix~\ref{sec:rh_close_scale}, does not enter this pooled denominator):
    % 2026-10-17 (user, E8): parked -- the bullet led with the closing-interval model. That display is now
    % Appendix~\ref{sec:rh_close_scale}. Was: written for the closing-interval model, whose rows t' are trading
    % days (i = c suppressed), on the estimation window W_t of Section sec:rh_methods_setup,
    % z^{(k)}_t = (bar u^{(k)}_t - med_{t' in W_t} bar u^{(k)}_{t'}) / max{IQR^{(k)}_t, 0.01 max_{t' <= t} IQR^{(k)}_{t'}, IQR^{*(k)}_t},
    % IQR^{(k)}_t = IQR{bar u^{(k)}_{t'} : t' in W_t}.
    % IQR^{*(k)}_t: IQR of bar u^{(k)}_{t'}, t' < t, over rows whose k-interval window is fully observed.
    % Mode at least 20% of rows: IQR* also excludes rows at the mode. IQR* = 1 until 8 such rows exist.
    % bar a^{(k)}, bar chi^{(k)} not standardized. HAR components and calendar dummies: IQR* = 0.
    % t < W: the first window's statistics.
    % 2026-10-16 (formula pass): parked -- "with W_t = {t - W, ..., t - 1}" repeated the window defined in Methods.
    % 2026-10-15 (consolidation): parked -- the window row was s (now t'; s is the periodicity factor):
    % z^{(k)}_{t} = (bar u^{(k)}_{t} - med_{s in W_t} bar u^{(k)}_{s}) / max{...}.
      \[
        z^{(k)}_{t,i} =
        \begin{cases}
          \dfrac{\bar u^{(k)}_{t,i}
            - \operatorname{med}\{\bar u^{(k)}_{t',i'} : (t',i')\in\mathcal{W}_{t,i},\ a_{t',i'}=1\}}
            {\max\bigl\{\mathrm{IQR}^{(k)}_{t,i},\ 0.01\, M^{(k)}_{t,i}\bigr\}}
            & \text{if at least 512 such intervals},\\[8pt]
          0 & \text{otherwise},
        \end{cases}
      \]
      \[
        \mathrm{IQR}^{(k)}_{t,i}
          = \mathrm{IQR}\bigl\{\bar u^{(k)}_{t',i'} : (t',i')\in\mathcal{W}_{t,i},\ a_{t',i'}=1\bigr\}.
      \]
      \begin{itemize}
      \item On the initial estimation window, both statistics are computed using that window's own rows; an
        interquartile range below $10^{-12}$ is replaced by $1$ on the initial window and, at subsequent grid
        points, retains the previous scale factor.
      \item The indicator averages $\bar a^{(k)}$ and $\bar\chi^{(k)}$ are left unstandardized.
      \item The target HAR components and calendar variables do not use the availability-restricted denominator
        above; instead, their median and interquartile range are computed across all rows of
        % 2026-10-18 (user, new-notes 8): parked -- forward reference to the IQR* definition added. Was: "$\mathcal{W}_{t,i}$ with $\mathrm{IQR}^{*(k)} = 0$, subject to the same $0.01$ floor relative to the"
        $\mathcal{W}_{t,i}$ with $\mathrm{IQR}^{*(k)} = 0$ ($\mathrm{IQR}^{*(k)}$ is defined in Appendix~\ref{sec:rh_close_scale}), subject to the same $0.01$ floor relative to the
        running maximum interquartile range of that column.
      \end{itemize}
    % 2026-10-17 (user, E8): parked -- the one-sentence intraday rule, and the note. The denominator is the
    % display above. Was: median and IQR over rows with a_{t',i'}=1; fewer than 512 such rows give z^{(k)}=0.
    % \editc{blue}{E8}{Decision: write the intraday denominator, or state that the display carries over. The display standardizes closing-interval rows $t'\in\mathcal W_t$ with denominator $\max\{\mathrm{IQR}^{(k)}_t,\ 0.01\max_{t'\le t}\mathrm{IQR}^{(k)}_{t'},\ \mathrm{IQR}^{*(k)}_t\}$, and sets $\mathrm{IQR}^{*(k)}_t=0$ on HAR components and calendar dummies. The intraday scheme, which is the one scored in the accuracy table, has one extra sentence: median and IQR over rows with $a_{t',i'}=1$, and $z^{(k)}=0$ when fewer than 512 such rows exist. That sentence does not say whether the 0.01 floor and the $\mathrm{IQR}^{*}=0$ exceptions apply to those rows.}
    % scaling.py:215-290 (IQR < 1e-12 -> 1); executor.py:96-202 (IQR^{*(k)}_t: expanding, 480-row grid, mode of
    % the ladder column); scaling.py:310-392 + executor.py:206-288 (masked scaling, pooled only; outline flag 7).
    % Per bar, the availability erosion counts k in design rows (sessions) while the ladder counts bars
    % (executor.py:178-188; outline flag 8, unmeasured).
  \end{enumerate}
\item \textbf{Column screening.} At each estimation origin, we screen out constant and duplicate predictors over
  the active estimation window $\mathcal{W}_t$ by retaining the index set
  \[
    \mathcal{J}_t = \bigl\{ j :\ \bm{x}_{\cdot j} \text{ not constant},\ \bm{x}_{\cdot j} \ne \bm{x}_{\cdot j'}\
      \forall\, j' < j,\ j' \in \mathcal{J}_t \bigr\}, \qquad
    \bm{x}_{\cdot j} = (x_{t'j})_{t' \in \mathcal{W}_t} .
  \]
  % 2026-10-15 (consolidation): parked -- window row s -> t'. Was: \bm{x}_{\cdot j} = (x_{sj})_{s \in \mathcal{W}_t}.
  \begin{itemize}
  \item Specifically, column equality is evaluated bitwise, and whenever two columns are identical on
    $\mathcal{W}_t$, the first column in order is retained.
  % 2026-10-17 (user): parked -- closing-interval material (moved to Appendix app:rh_closing_predictors). Was:
  % \item Under the closing-interval specification with the full predictor set, this screening procedure retains a
  %   median of \cmKeptAll{} out of \cmPNewAll{} candidate columns across rolling windows.
  \end{itemize}
  % window_mask.py:27-46 (float64 bytes).
% 2026-10-15 (user): parked -- Table 3 (tab:rh_dm_design) is parked. Was:
% \item \textbf{Predictor sets} (Table~\ref{tab:rh_dm_design})\textbf{:} the HAR benchmark has \cmPNewBase{} columns;
% 2026-10-17 (user): parked -- closing-interval material (closing column counts moved to Appendix app:rh_closing_predictors). Was:
% \item \textbf{Predictor set dimensions.} In the closing-interval specification, the HAR benchmark contains
%   \cmPNewBase{} candidate columns, whereas the full predictor set contains \cmPNewAll{} columns:
%   \[
%     \cmPNewAll = \underbrace{\cmPNewBase}_{\bar y^{(k)},\ \bm D}
%       + \underbrace{\rhfiNExog \times \rhdmNHorizonsClose}_{\bar u^{(k)}}
%       + \underbrace{\rhfiNExog \times \rhdmNHorizonsClose}_{\bar a^{(k)}}
%       + \underbrace{\rhdmNActive \times \rhdmNHorizonsClose}_{\bar\chi^{(k)}}, \qquad
%     \cmPNewBase = \rhdmNHorizonsClose + \rhdmNCalendar .
%   \]
\item \textbf{Predictor set dimensions.} The column counts of the closing-interval specification are given in
  Appendix~\ref{app:rh_closing_predictors}.

  \begin{itemize}
  \item In the pooled intraday specification, the full predictor set comprises \rhdmNColsPooled{} candidate
    columns, constructed from \rhdmNSeriesPooled{} exogenous series, \rhdmNHorizonsPooled{} horizons, and
    \rhdmNEdgePooled{} interaction terms between the target HAR components and the opening- and
    closing-interval dummy variables. Note that these column counts still reflect the
    % 2026-10-17 (user): parked -- closing-interval material. Was: "Note that these column counts, as well as $\cmPNewBase$, still reflect the"
    calendar columns and interaction terms discussed in note E9 above:
    \[
      \rhdmNColsPooled = \underbrace{\rhdmNHorizonsPooled + \rhdmNCalendar}_{\bar y^{(k)},\ \bm D}
        + \underbrace{\rhdmNEdgePooled}_{\bar y^{(k)} \times \text{open, close}}
        + \rhdmNSeriesPooled \times \rhdmNHorizonsPooled + \rhdmNSeriesPooled \times \rhdmNHorizonsPooled
        + \rhdmNActivePooled \times \rhdmNHorizonsPooled .
    \]
    % 2026-10-19 (user): parked -- the pooled model includes the FOMC series (author decision). Was: "Here, the $\rhdmNSeriesPooled$ intraday exogenous series encompass every category in Section~\ref{sec:rh_data} except the FOMC group (which accounts for the remaining $\bkSizeFomc$ series)."
    Here, the $\rhdmNSeriesPooled$ intraday exogenous series encompass every category in
    Section~\ref{sec:rh_data}, the FOMC group ($\bkSizeFomc$ series) included.
    Appendix Table~\ref{tab:rh_dm_dictionary} lists each individual series.
    % 2026-10-19 (user): base 5 across the paper -- the column counts above are recomputed on the base-5 pooled panel.
    % 2026-10-19 (user): parked -- the pooled model includes the FOMC series (author decision); E35 recounted on the panel with FOMC. Was:
    %     \editc{orange}{E35}{Decision: base 5 across the paper (user, 2026-10-19).\\Problem: $\rhdmNColsPooled$ and the formula above now count the base-5 pooled panel built on 19 October 2026, before the cuts of notes E6 and E9; no printed forecast was fitted on that panel.\\
    % The earlier pooled trees and the intraday OLS were fitted on the base-2 panel, 1,144 columns from 12 horizons (note E34). The base-5 panel (experiments/pooled\_trees\_base5.py, design stage: the repo's own preparation with the data fixes of 15 October 2026 and the causal mode rule) has $\rhdmNHorizonsPooled$ horizons and $\rhdmNEdgePooled$ interaction columns. Under the causal rule $\rhdmNActivePooled$ series carry an activity indicator, the number of observations and the three StockTwits series, where the old full-sample rule gave 5, the number of observations and the four volatility-demand series (note S4a). The design of the base-5 tree refits also drops the activity averages (E6) and the calendar after the overnight indicator together with the interaction columns (E9): 529 columns on all $\rhdmNIntervals$ intervals and 528 on the 13 regular-hours bars, where the overnight indicator is dropped too. Those two counts have no macro yet.}
    \editc{orange}{E35}{Decision: base 5 across the paper, and the pooled model includes the FOMC series (user, 2026-10-19).\\Problem: $\rhdmNColsPooled$ and the formula above now count the base-5 pooled panel with the FOMC series, built on 19 October 2026, before the cuts of notes E6 and E9; no printed forecast was fitted on that panel.\\
The earlier pooled trees and the intraday OLS were fitted on the base-2 panel without the FOMC series, 1,144 columns from 12 horizons and 43 series (note E34). The base-5 panel (experiments/pooled\_trees\_base5.py, design stage: the repo's own preparation with the data fixes of 15 October 2026, the FOMC calendar fix and the causal mode rule) has $\rhdmNHorizonsPooled$ horizons, $\rhdmNSeriesPooled$ series and $\rhdmNEdgePooled$ interaction columns. Under the causal rule $\rhdmNActivePooled$ series carry an activity indicator, the number of observations, the three StockTwits series and the FOMC statement-interval and statement-day dummies, where the old full-sample rule gave 5 on the panel without FOMC, the number of observations and the four volatility-demand series (note S4a). With the activity averages dropped (E6) and the calendar cut after the overnight indicator together with the interaction columns (E9), the design has 577 columns on all $\rhdmNIntervals$ intervals and 576 on the 13 regular-hours bars, where the overnight indicator is dropped too. Those two counts have no macro yet.}
    % 2026-10-17 (user, E5): parked -- the series are listed in the appendix table, and the FOMC group is the
    % set left out of the intraday matrix. Was:
    % \editc{blue}{E5}{Decision: list the series in the intraday matrix.\\Problem: the $\rhdmNColsPooled$ intraday columns are built from $\rhdmNSeriesPooled$ of the $\rhfiNExog$ series, and neither the series that are in nor the series that are left out are named.\\
    % The closing-interval formula adds one block of columns for each exogenous series, and it says there are $\rhfiNExog$ such series. The names of those series are not written in this section, and the dictionary table that would list them is commented out of the floats. The intraday formula uses a smaller count, $\rhdmNSeriesPooled$ series. The series that make up the difference, $\rhfiNExog-\rhdmNSeriesPooled$ of them, are not named either. Without that list the $\rhdmNColsPooled$ intraday columns cannot be written out from the paper.}
    % 2026-10-15 (finish, phase 2): typed 6, 16, 7, 1,144, 43, 12, 24, 5 -> \rhdmNHorizonsClose, \rhdmNCalendar,
    % \rhdmNActive, \rhdmNColsPooled, \rhdmNSeriesPooled, \rhdmNHorizonsPooled, \rhdmNEdgePooled, \rhdmNActivePooled.
    % Pooled panel of record is pre-FOMC (43 series; 5 zero-inflated: numobs + 4 voldemand). Current code
    % builds 1,264 = 52 + 47*12*2 + 7*12 (outline flag 10). Session-edge columns: executor.py:454-459.
  \end{itemize}
\item \textbf{Causal construction and absence of look-ahead.} Every rolling statistic indexed by $(t,i)$ is
  % 2026-10-17 (user): parked -- closing-interval material. Was: "ensuring that all closing-interval predictors are in the information set at the 15:30 ET forecast origin."
  computed strictly from intervals prior to $(t,i)$, ensuring that every predictor of interval $(t,i)$ is in the
  information set at its forecast origin (15:30 ET for the scored closing interval $i = c$). This includes the HAR trailing averages, the diurnal
  factor $s^2_{t,i}$, the winsorization quantiles, the median baseline $\rho_{t,i}$, the standardization moments,
  and the active column set $\mathcal{J}_t$.
  \begin{itemize}
  \item We disclose three full-sample quantities used in preprocessing: the full-sample mode used to define
    $\chi_{t,i}$, the within-interval standard-deviation floor $\bar\sigma^{(u)}$, and the rule assigning series
    to the inverse hyperbolic sine transformation.
    % 2026-10-18 (user, second-list N4): rerun-needed flag (old-rule disclosure).
    \editc{orange}{S4b}{Decision: rerun needed under the causal rule (note S4a).\\Problem: the full-sample mode is the old rule, used by the printed forecasts of 6--7 October 2026. Under the causal rule in src/backtest/executor.py the mode is decided from earlier prints only, so it is no longer a full-sample quantity; drop it from this list after the rerun.}
  \item As noted in Section~\ref{sec:rh_data}, the FOMC proximity variables rely solely on the pre-announced
    meeting schedule, which is publicly known ahead of each interval.
  \item Any missing interval prior to a series' initial observation is filled with that first observed print
    (step 4).
  \item Finally, standardization statistics on the first estimation window are computed exclusively from
    burn-in observations.
  % 2026-10-16 (user): the imputation median left this list. It is the expanding median of earlier
  % observed values. Parked -- "the imputation median" as a full-sample element.
  \end{itemize}
  % The previous text said all full-sample elements touch warm-up rows only. False: the neutral median fills
  % every a_{t,i} = 0 row, including the 2023-24 dead tails inside the test period. Also full-sample: the FOMC
  % dead-tail cut (outline section 4, "Full-sample elements").
\end{itemize}

